% ============================================================
%  SHPN-NESM-Notes.tex
%  Working notes from scientific discussion — July 2026
%  Topics: PreemptionCheck/trajectories, POSet as entropic ratchet,
%          ATP bath concern, kBT origin, reversibility in ST,
%          N0 sweep and Fujita calibration defence
% ============================================================

\documentclass[11pt,a4paper]{article}
\usepackage[utf8]{inputenc}
\usepackage[T1]{fontenc}
\usepackage[english]{babel}
\usepackage{amsmath,amssymb}
\usepackage{booktabs}
\usepackage{geometry}
\usepackage{hyperref}
\usepackage{microtype}
\usepackage[font=small,labelfont=bf]{caption}
\geometry{margin=2.5cm}
\hypersetup{colorlinks=true,linkcolor=blue!50!black,citecolor=blue!50!black}

\newcommand{\kB}{k_{\mathrm{B}}}
\newcommand{\EPR}{\dot{\sigma}}

\title{\textbf{SHPN-NESM Notes}\\
  \large Working notes from scientific discussion --- July 2026}
\author{Eug\^{e}nio Sim\~{a}o}
\date{2026-07-10}

\begin{document}
\maketitle
\tableofcontents
\newpage

% ─────────────────────────────────────────────────────────────────────────────
\section*{Note 1 — Reversibility in Stochastic Thermodynamics}
% ─────────────────────────────────────────────────────────────────────────────

Reversibility is \emph{central} to stochastic thermodynamics, not an
exception. A process is reversible iff $P[\omega] = P[\tilde\omega]$ for
every trajectory, which implies $\sigma[\omega] = 0$ (detailed balance at
path level). The Crooks fluctuation theorem $P_F[\omega]/P_R[\tilde\omega]
= e^{\sigma[\omega]}$ quantifies \emph{how far} a trajectory deviates from
reversibility. At $\sigma = 0$ (ratio $= 1$) forward and reverse are
equally probable. The three regimes: equilibrium ($\sigma = 0$ always),
near-equilibrium ($\sigma \ll 1$, Onsager linear response, FTs testable),
far-from-equilibrium ($\sigma \gg 1$, FTs exact but inaccessible, NESS
regime). Sporulation sits firmly in the third regime. ST does not
\emph{require} reversibility; it \emph{measures} the degree of
irreversibility. $\dot\sigma \geq 0$ is the Second Law at trajectory
level, valid for any system.

% ─────────────────────────────────────────────────────────────────────────────
\section*{Note 2 — Origin of $k_BT$ in a Single-Cell NESM System}
% ─────────────────────────────────────────────────────────────────────────────

$k_BT$ does not arise from the cell's internal temperature; it enters
because the cell is coupled to an \emph{external thermal bath} — the
aqueous cytoplasm in thermal equilibrium with the surrounding medium at
$T \approx 310\,\text{K}$. Local detailed balance:
$k^+_\rho / k^-_\rho = \exp(\Delta\mu_\rho / k_BT)$ encodes the bath
temperature, not a cellular property. Every molecular collision and ATP
hydrolysis event occurs in aqueous solution at $T = 310\,\text{K}$;
the solvent is what makes reactions stochastic. The cell is
\emph{chemically} non-equilibrium (sustained by ATP/nutrient gradients)
but \emph{thermally} equilibrated (coupled to a large solvent bath with
fast relaxation time $\tau_{\rm relax} \sim \text{ps}$, much faster than
biochemical timescales). This separation is precisely what makes
stochastic thermodynamics applicable at the single-cell level. $W/k_BT
= 3.72\times10^7$ means commitment dissipates that many thermal energy
quanta — not enough thermal noise exists in the bath to drive the reverse
process spontaneously.

% ─────────────────────────────────────────────────────────────────────────────
\section*{Note 3 — PreemptionCheck, Cascade EPR Silence, and $\gamma \neq$ MF}
% ─────────────────────────────────────────────────────────────────────────────

PreemptionCheck modifies the effective propensity at trajectory level:
$\EPR_t(t) = 0$ when PreemptionCheck is false, $\Phi_t \ln(\Phi_t^+/
\Phi_t^-)$ when true. This creates \emph{structured EPR silence intervals}
in the trajectory — periods when the entire cascade is below threshold and
produces no stochastic EPR. Because PreemptionCheck couples all Layer-A
transitions through the signal hierarchy $\lambda$, when it fails for one
transition it fails for all downstream transitions simultaneously.
Consequence: EPR contributions from different transitions are NOT
independent Poisson events (as in a standard CRN). They share a common
on/off switch. Variance of total EPR is suppressed relative to the
Poisson/mean-field prediction.

This is the physical origin of $\gamma_{\rm sub} = 0.150 < 0.5$: near
$N_c$, the cascade turns on and off collectively (joint EPR bursts),
not transition-by-transition (independent Poisson). The variance
suppression is structural — enforced by the cascade topology — not a
statistical artefact.

Three mechanisms contribute:
\begin{enumerate}
  \item \textbf{Sequential amplification} (KinA $\to$ Spo0F $\to$ Spo0A):
    fluctuation in KinA propagates multiplicatively through the chain.
  \item \textbf{Nonlinear inhibitor gate} (SinR $< 7.5\,\mu$M):
    sharp threshold creates history-dependent gating.
  \item \textbf{PreemptionCheck coupling}: Layer-A transitions are jointly
    suppressed when the cascade is below threshold.
\end{enumerate}

Critical slowing down signature: near $N_c$, the marking of SinR
fluctuates around $\theta_{\rm eff} = 7.5\,\mu$M. The cascade turns on
and off rapidly (EPR flickering) before finally committing. Mean duration
of active intervals increases as $N_0 \to N_c$. The total
$\sigma[\omega] = \int \EPR(t)\,dt$ is the integral over active intervals
only.

Standard CRN (no hierarchy) predicts $\mathrm{Var}[\sigma[\omega]] \sim
\langle\sigma[\omega]\rangle$ (Poisson / MF). With PreemptionCheck:
$\mathrm{Var}/\langle\sigma\rangle < 1$ near $N_c$. This ratio is the
thermodynamic fingerprint of the non-MF universality class, enforced by
cascade topology.

% ─────────────────────────────────────────────────────────────────────────────
\section*{Note 4 — POSet $\lambda$ as an Entropic Ratchet}
% ─────────────────────────────────────────────────────────────────────────────

\textbf{Core insight:} The POSet layer assignment encodes a directed
state space where each level transition is irreversible because the free
energy cost of ascending is consumed from the chemical bath and cannot
be recovered.

\textbf{Formal statement:}
$$\frac{W(0 \to n)}{W(n \to 0)} = \prod_{k=0}^{n-1} \frac{k^+_k}{k^-_k}
= \exp\!\left(\frac{\sum_k \Delta\mu_k}{k_BT}\right)
= e^{\mathcal{A}_{\rm total}/k_BT}$$

For $\mathcal{A}_{\rm total} = 3.72\times10^7$: forward-to-reverse ratio
$= e^{3.72\times10^7} \approx 10^{16{,}000{,}000}$.

\textbf{Cascade depth scaling:}
$$\frac{W}{\kB T} \sim n \cdot \bar{\mathcal{A}} \cdot \bar{n}_{\rm firings}$$
Deeper POSet $\to$ more levels $\to$ larger $W/k_BT$ $\to$ stronger
irreversibility. Falsifiable: single-level circuits ($n=1$) should have
$W/k_BT \sim 10^3$; three-level circuits ($n=3$) $\sim 10^7$;
both should be \emph{structurally} irreversible beyond $n = 3$ for any
realistic $\Delta\mu$.

\textbf{Why PreemptionCheck reinforces this:}
The reverse trajectory $\tilde\omega$ requires $c$ to fire before $b$,
$b$ to fire before $a$. But PreemptionCheck sets $\Phi_b(M_{\rm before\,a})
= 0$. Therefore $P[\tilde\omega] = 0$ under forward dynamics. The POSet
converts a \emph{probabilistically suppressed} reverse path into a
\emph{structurally forbidden} one.

\textbf{Key sentence for manuscript:}
\begin{quote}
\emph{The POSet layer assignment $\lambda$ encodes an entropic ratchet:
each level of the cascade consumes a fixed free energy quantum from the
chemical bath, PreemptionCheck prevents descent without repaying that
quantum, and the accumulated affinity $\mathcal{A}_{\rm total} =
\sum_k \Delta\mu_k \gg k_BT$ makes the time-reversed trajectory
thermodynamically inaccessible. Biological commitment irreversibility is
not an emergent statistical property --- it is structurally enforced by
the cascade architecture.}
\end{quote}

\textbf{Biological parallel:}
Shallow circuits ($n = 1$, e.g.\ direct TF binding) have modest $W/k_BT$
and could in principle be reversed by external perturbation. Deep circuits
($n \geq 3$, phosphorelay + inhibitor gate) are effectively irreversible.
The transition from reversible to irreversible commitment is quantifiable
as a cascade depth threshold.

% ─────────────────────────────────────────────────────────────────────────────
\section*{Note 5 — ATP Bath Concern: Physicist's Challenge and Response}
% ─────────────────────────────────────────────────────────────────────────────

\textbf{Challenge:} The v9 model treats ATP as a dedicated bath for the
sporulation subnetwork. In reality other processes (translation $\sim 80\%$
of ATP demand, replication, membrane synthesis) draw from the same pool
concurrently. This means $\sigma_{\rm interface} \neq 0$ and the simple
subsystem decomposition $\sigma_{\rm total} = \sigma_{\rm sub} +
\sigma_{\rm rest} + \sigma_{\rm interface}$ is invalid.

\textbf{Scale separation defence:}
\begin{itemize}
  \item The \emph{signalling} layer (KinA cycling) consumes $\sim 8\,\mu$M
    ATP over 6\,h — negligible vs.\ 5000\,\mu$M pool.
  \item During starvation (the commitment window): ribosomes hibernate,
    replication halts, flagella downregulate. Other ATP drains are
    suppressed precisely when sporulation is active.
  \item The qualitative EPR decomposition (KinA 97.5\% of stochastic EPR)
    is a \emph{relative} structural property — robust to absolute pool
    level changes.
\end{itemize}

\textbf{N0 sweep partial answer:}
Sweeping INITIAL\_NUTRIENTS from 100 to 2200\,$\mu$M covers the full
range from quiescent starvation (other processes shut down, ATP
near-dedicated to sporulation) to abundance (sporulation suppressed,
question moot). The critical regime near $N_c = 1346\,\mu$M is where
competition is most relevant — and also where the most important results
originate.

However: the sweep varies the \emph{trigger condition} (nutrient
availability), not the \emph{competition intensity} (concurrent ATP
drain). These are orthogonal axes. The sweep does not fully address the
concern.

\textbf{Fujita calibration argument (strongest response):}
The model's $N_c = 1346\,\mu$M and sporulation fractions were
implicitly calibrated against real biological experiments (Fujita \&
Losick 2005) in which ALL concurrent ATP consumers were active. The
model reproduces the experimental 38\% sporulation at CH conditions.
Therefore the structural results ($N_c$, $\gamma_{\rm sub}$, EPR
decomposition ratios) are empirically grounded in data that already
includes competition — even if the mechanism is not explicitly modelled.

\textbf{What remains genuinely sensitive:}
Absolute values of $W/k_BT$ and EPR (kBT/s) would shift if background
ATP consumption were explicitly included. These should be interpreted as
``within the sporulation subnetwork model under starvation conditions,''
not as total cellular thermodynamics.

\textbf{Next model step:}
Add explicit ATP regeneration transitions + background consumption sink
scaling with $N_0$ $\to$ dynamic ATP NESS $\to$ rigorous Markovian
reservoir for the signalling layer.

% ─────────────────────────────────────────────────────────────────────────────
\section*{Note 6 — PreemptionCheck as Artifact vs.\ Thermodynamic Emergence}
% ─────────────────────────────────────────────────────────────────────────────

\textbf{The theoretical critique:}
In canonical stochastic thermodynamics (Seifert, Schnakenberg) every
reaction has a non-zero reverse rate and irreversibility emerges from
non-zero affinities $\Delta\mu > 0$.
PreemptionCheck sets the propensity to \emph{exactly zero} when the
upstream cascade is inactive.
A theorist argues: this is an architectural constraint, not a thermodynamic
consequence.
When $P[\tilde\omega] = 0$ exactly, the Crooks theorem is trivially
satisfied (denominator zero) and $\sigma[\omega] \to \infty$ — not a
meaningful thermodynamic statement but a tautology from the imposed
disconnection.

\textbf{Where the critique is valid:}
PreemptionCheck IS an architectural constraint. A fully thermodynamically
grounded model would have non-zero reverse propensities for every
transition (however small), with $W/k_BT$ large because of affinities,
not because reverse paths are set to zero.
The $W/k_BT = 3.72 \times 10^7$ result is partly an idealisation of
strong kinetic suppression — the physical reverse rate is very small but
not exactly zero.

\textbf{Where the critique overreaches:}
PreemptionCheck encodes a \emph{biochemical fact}, not a modelling
convenience:
\begin{enumerate}
  \item The phosphorelay is causally ordered by \emph{molecular recognition}:
    Spo0F is phosphorylated by KinA through direct protein-protein contact.
    Without KinA~P present, the reaction has no substrate — not just
    thermodynamically suppressed, but biochemically impossible.
  \item The reverse pathway (dephosphorylation) uses \emph{different enzymes}
    (RapA, RapB, Spo0E phosphatases) with different rate constants.
    It is not the time-reverse of the forward step — it is a
    different molecular process.
  \item PreemptionCheck is the idealised limit of $k^-_k \to 0$ as
    protein specificity becomes perfect.
    In the real system the reverse rate is very small; zero is the
    mathematical idealisation of that limit.
\end{enumerate}

\textbf{Two independent irreversibility mechanisms:}
\begin{center}
\begin{tabular}{lll}
\hline
Mechanism & Origin & SHPN encoding \\
\hline
Thermodynamic driving & $\Delta\mu_k \gg k_BT$ (ATP hydrolysis) &
  Rate functions, local detailed balance \\
Biochemical causal ordering & Protein-protein specificity &
  PreemptionCheck ($\Phi = 0$ when upstream inactive) \\
\hline
\end{tabular}
\end{center}

Neither alone is sufficient. Together they enforce irreversibility at
both the energetic level (affinity $\mathcal{A}_{\rm total} \gg k_BT$)
and the kinetic level (reverse pathway is a different process, not
simply thermodynamically suppressed).

\textbf{Consequence for the FWD perturbation result (NET-T3):}
FWD aborts commitment ($38\% \to 0\%$) because the energetics are
overridable by external perturbation.
But you cannot restart the reverse cascade by adding nutrients —
the biochemical specificity constraint is not energetically reversible.
This is observable in the asymmetry: FWD completely abolishes but REV
only weakly triggers ($0\% \to 2\%$).
The 2\% represents the rare stochastic events where the cascade
completes despite the reversed nutrient step; the near-zero rate is
the thermodynamic suppression.
The 0\% for FWD represents structural abortion (cascade already at
high affinity toward vegetative basin, irreversible at the biochemical
specificity level).

\textbf{Manuscript statement (Theory II §4.3):}
\begin{quote}
\emph{Biology achieves commitment irreversibility through two independent
mechanisms: thermodynamic driving ($\Delta\mu_k \gg k_BT$ at each
ATP-consuming step) and biochemical causal ordering (protein specificity,
encoded by PreemptionCheck). Biological commitment irreversibility is
therefore doubly enforced: thermodynamically by the accumulated affinity
of the cascade, and biochemically by the protein-specificity constraints
that make the reverse pathway a different molecular process entirely.}
\end{quote}

% ─────────────────────────────────────────────────────────────────────────────
\section*{Note 7 — How to Categorise NESM + SHPN as a Scientific Contribution}
% ─────────────────────────────────────────────────────────────────────────────

\textbf{Date:} 2026-07-13.
The contribution sits at four levels simultaneously.

\subsection*{Level 1 — A formal language (Computer Science / Theoretical Biology)}

SHPN is a mathematically defined formal object: a 13-tuple with a provable
firing semantics, typed arc grammar, and layer assignment $\lambda$.
This places it in the tradition of Chemical Reaction Network Theory
(Feinberg 1972), Stochastic Petri Nets (Molloy 1982), and Harel Statecharts
(1987), but with a richer type system that handles signal hierarchy and
hybrid stochastic/ODE dynamics.
\begin{quote}
\emph{Contribution: a new formal language for regulatory circuits with
provable compositional properties.}
\end{quote}

\subsection*{Level 2 — A computational substrate for stochastic thermodynamics (Physics)}

By mapping SHPN arc semantics onto thermodynamic ports (normal arcs =
stoichiometric flows; signal-flow = regulatory dissipation channels;
test/inhibitor = zero-EPR sensing), the formalism becomes a natural
substrate through which the machinery of NESM (Seifert 2012,
Rao-Esposito 2018) operates.
The $\tau$-leaping engine is a Master equation approximation; the five
thermodynamic quantities emerge from topology and dynamics without
additional assumptions.
\begin{quote}
\emph{Contribution: a natural computational substrate for stochastic
thermodynamics of regulatory decision circuits.}
\end{quote}

\subsection*{Level 3 — A measurement methodology (Systems Biology / Biophysics)}

The five quantities ($\eta_L$, $W/k_BT$, $N_c$, $\gamma_{\rm sub}$,
EPR decomposition) are computable from any validated SHPN model without
new experiments.
This provides a new measurement layer on top of existing biological models
--- extracting thermodynamic information that was always implicitly present
in the network topology but previously inaccessible.
\begin{quote}
\emph{Contribution: a methodology for computing the thermodynamic
fingerprint of any topology-known commitment circuit.}
\end{quote}

\subsection*{Level 4 — A conceptual unification (Interdisciplinary)}

\begin{center}
\begin{tabular}{p{4cm}p{4.5cm}p{4.5cm}}
\toprule
Tradition & What NESM+SHPN inherits & What it adds \\
\midrule
NESM / stochastic thermodynamics & EPR, fluctuation theorems, NESS &
  Biological regulatory hierarchy, hybrid dynamics \\
Petri net formalism & Firing semantics, compositionality &
  Thermodynamic arc types, POSet irreversibility \\
Systems biology (ODE) & Biochemical rate functions &
  Single-cell thermodynamics, stochastic decision layer \\
Single-cell biology & Bet-hedging, commitment &
  Quantitative thermodynamic characterisation \\
\bottomrule
\end{tabular}
\end{center}

\subsection*{Most precise academic categorisation}

\begin{quote}
\emph{A formal computational framework that provides stochastic
thermodynamics as a derived property of regulatory circuit topology,
applicable to any cell-fate commitment system with known network
architecture.}
\end{quote}

This places it in the category of \textbf{theory-enabling frameworks}:
like Gillespie's SSA (1977) enabled stochastic kinetics simulation, or
CRNT (Feinberg 1972) enabled stability analysis of biochemical networks,
NESM+SHPN enables thermodynamic characterisation of regulatory circuits.

\subsection*{What it is NOT}

\begin{itemize}
  \item Not a theory of cell biology (does not explain why circuits are
    wired as they are).
  \item Not a simulation tool (the SHPN engine is the tool; the framework
    is the language + quantities).
  \item Not a phenomenological model (quantities are derived, not fitted).
\end{itemize}

\subsection*{Journal framing guide}

\begin{center}
\begin{tabular}{p{3.5cm}p{3.5cm}p{5.5cm}}
\toprule
Target & Lead category & Framing \\
\midrule
Physical Review X / PRE & Computational substrate for NESM &
  ``First formal substrate through which stochastic thermodynamics is
  computable from biological regulatory topology'' \\
eLife / PLOS CB & Measurement methodology &
  ``Five topology-derivable quantities characterising the thermodynamic
  cost of cell-fate commitment'' \\
Cell Systems / Nature Methods & Formal language + tool &
  ``Hybrid stochastic/ODE framework in which thermodynamic properties
  emerge from network architecture'' \\
Biophysical Journal & Unification &
  ``NESM and regulatory network formalism unified, validated on
  B.\ subtilis'' \\
\bottomrule
\end{tabular}
\end{center}

% ─────────────────────────────────────────────────────────────────────────────
\section*{Note 8 — Cell Dimensions, Reaction Counts, and Genome-to-Pathway Inference}
% ─────────────────────────────────────────────────────────────────────────────

\textbf{Date:} 2026-07-15.

\subsection*{Single-cell dimensions (micron scale)}

\begin{center}
\begin{tabular}{llll}
\toprule
Cell type & Dimensions & Volume & Notes \\
\midrule
\textit{B.~subtilis} (sporulating) & 2--4\,$\mu$m $\times$ 0.8--1.0\,$\mu$m & $\sim1$--2\,fL & Our model \\
\textit{E.~coli} & 1--2\,$\mu$m $\times$ 0.5\,$\mu$m & $\sim1$\,fL & \\
\textit{S.~cerevisiae} & 5--10\,$\mu$m diameter & $\sim100$\,fL & \\
Mammalian cell & 10--30\,$\mu$m diameter & $\sim1$--10\,pL & \\
\bottomrule
\end{tabular}
\end{center}

Molecule count at 1\,$\mu$M in a 1\,fL cell:
$N = C \times N_A \times V = 1\,\mu\text{M} \times 6.022\times10^{23}
\times 10^{-15}\,\text{L} \approx 600$ molecules.
This confirms the stochastic regime for KinA$\sim$P ($\sim$1.5\,$\mu$M
$\approx$ 900 molecules): Poisson fluctuations $\sim 3\%$ of mean,
sufficient to drive bet-hedging.

\subsection*{Number of reactions in a single cell}

\begin{center}
\begin{tabular}{lrrl}
\toprule
Organism & Metabolic reactions & Protein-coding genes & Ratio \\
\midrule
\textit{E.~coli} & $\sim$2,077 & $\sim$4,400 & 0.47 \\
\textit{B.~subtilis} & $\sim$1,200 & $\sim$4,100 & 0.29 \\
\textit{S.~cerevisiae} & $\sim$3,500 & $\sim$6,000 & 0.58 \\
Human cell & $\sim$13,000 & $\sim$20,000 & 0.65 \\
\bottomrule
\end{tabular}
\end{center}

The v9 SHPN model has 36 transitions: $\sim$3\% of \textit{B.~subtilis}
metabolic reactions.
During sporulation starvation, $\sim$700 reactions shut down (translation
reduced 10$\times$, replication halted), leaving $\sim$500 active ---
so the sporulation subnetwork represents $\sim$7\% of active chemistry.
This contextualises the Markovian reservoir assumption: the bath contains
$\sim$500 concurrent reactions, not the full 1,200.

\subsection*{Can biochemical pathways be discovered from genome data?}

\textbf{Yes, routinely --- with important caveats.}

\textbf{What genome sequence gives directly:}
\begin{itemize}
  \item Enzyme identification by homology (BLAST/HMMer vs.\ KEGG, UniProt,
    MetaCyc): $>40\%$ amino acid identity $\to$ function inferred.
  \item Metabolic network reconstruction (PathwayTools, ModelSEED,
    CarveMe): automated assembly of reactions from annotated genes.
  \item Operon structure: gene proximity $\to$ pathway co-membership
    (bacteria).
\end{itemize}

\textbf{What cannot be inferred from sequence alone:}
\begin{center}
\begin{tabular}{ll}
\toprule
Problem & Why it fails \\
\midrule
Novel enzymes & No homology $\to$ unknown function ($\sim$30--40\% of
  any new genome) \\
Enzyme promiscuity & One enzyme may catalyse multiple reactions \\
Regulatory wiring & TF$\to$gene control requires ChIP-seq/expression data \\
Kinetic parameters & $K_m$, $V_{\max}$, rate constants require measurement \\
Post-translational control & Allosteric, phosphorylation effects invisible \\
\bottomrule
\end{tabular}
\end{center}

\textbf{Modern pipeline:}
Genome $\to$ gene calling $\to$ homology search $\to$ EC numbers $\to$
metabolic reconstruction $\to$ gap filling $\to$ FBA predictions $\to$
experimental validation $\to$ curated model.

\textbf{Relevance to SHPN:}
The v9 model was built from decades of targeted biochemical experiments
(knockouts, reporter fusions, Western blots) --- not genome inference.
Genome-scale reconstruction gives the reaction inventory; it does not
give the causal ordering ($\lambda$ layers), inhibitor thresholds
($\theta_{\rm eff}$), or stochastic firing rates.
These require quantitative single-cell experiments (Fujita \& Losick
2005).

\textbf{Future opportunity:}
Once a pathway is genome-reconstructed and kinetically validated, the
NESM+SHPN framework provides its thermodynamic characterisation
essentially for free from the topology --- no additional experiments
needed.

% ─────────────────────────────────────────────────────────────────────────────
\section*{Note 9 — $G_E$ and $G_s$ as Markov Process Structures}
% ─────────────────────────────────────────────────────────────────────────────

\textbf{Date:} 2026-07-16.

\subsection*{$G_E$ as the Markov generator}

$G_E = (P, T, F)$ defines the stoichiometric structure of an underlying
continuous-time Markov chain (CTMC).
Each transition $t \in T$ contributes to the Master equation:
\begin{equation}
\frac{dP(M,t)}{dt} = \sum_{t \in T}
  \left[\Phi_t(M-\Delta_t)\,P(M-\Delta_t,t)
       - \Phi_t(M)\,P(M,t)\right]
\end{equation}
where $\Delta_t$ is the stoichiometric change vector and $\Phi_t(M)$ is
the propensity at state $M$.
\textbf{$G_E$ is the generator matrix of the Markov process:} it defines
which states are reachable from which and at what rates.

\subsection*{$G_s$ as a conditional enabling filter}

$G_s = (\Psi, F_s)$ does not define new Markov states.
It defines a \emph{conditional structure} on the generator via
PreemptionCheck:
\begin{equation}
\Phi_t^{\mathrm{eff}}(M) = \Phi_t(M)
  \cdot \mathbf{1}[\mathrm{PreemptionCheck}(t,M)]
\end{equation}
where the indicator depends on whether signal places $p_s \in \Psi$
upstream of $t$ in $G_s$ are above $\theta_{\mathrm{eff}}$.
In probability theory terms, $G_s$ defines a \textbf{filtration} ---
a coarsening of the state information.
Let $M_s = \{M(p) : p \in \Psi\}$ be the signal-place substate.
The combined $G_E \times G_s$ process is the $G_E$-process
\emph{conditioned} on the $G_s$ state.

\subsection*{Product structure and consequences}

\begin{center}
\begin{tabular}{p{3cm}p{5cm}p{5cm}}
\toprule
Concept & $G_E$ & $G_s$ \\
\midrule
Markov role & Generator matrix (stoichiometry + rates) &
  Conditional enabling filter (PreemptionCheck) \\
State space & Full marking $M: P \to \mathbb{N}_0$ &
  Signal substate $M_s: \Psi \to \mathbb{N}_0$ \\
EPR contribution & Mass flows through arcs &
  Threshold enforcement; zero EPR for pure gating \\
Stationarity & NESS or equilibrium of bare CRN &
  Creates absorbing subspaces \\
Causal role & Stoichiometric causality &
  Hierarchical causal ordering via $\lambda$ \\
\bottomrule
\end{tabular}
\end{center}

\subsection*{Absorbing subspaces — commitment as Markov absorption}

The most important Markov consequence of $G_s$ is the creation of
\textbf{absorbing subspaces}.
Before commitment: the system explores both attractor basins
(stochastic); $M(\text{SinR}) > 7.5\,\mu$M keeps T\_septation disabled.
Commitment event: stochastic SinI burst drops SinR below threshold;
T\_septation fires; forespore forms.
After commitment: the structural execution cascade (Layer B, ODE) runs
deterministically; no $G_s$-conditioned transition can reverse it.

In Markov terms: the committed state is an \textbf{absorbing class}.
The stationary distribution $P_{\mathrm{ss}}^{E \times s}(M)$ has support
only on the committed and vegetative attractors.

\textbf{The bet-hedging fraction is the absorption probability:}
$P(\mathrm{spor}) = P(\text{absorbed into committed class})$ from the
pre-commitment initial condition.
This is a classical Markov absorption problem --- the $N_0$ sweep
computes how this absorption probability changes with initial nutrient
marking.

\subsection*{EPR decomposition in $G_E \times G_s$ terms}

$G_s$ conditioning zeros propensities during PreemptionCheck-off
intervals (EPR silence, Note 3), suppressing both $G_E$ and $G_s$
contributions simultaneously:
\begin{equation}
\dot\sigma_{\mathrm{total}}
  = \underbrace{\sum_{t \in T_E} \Phi_t^{\mathrm{eff}}
      \ln \frac{\Phi_t^+}{\Phi_t^-}}_{\text{EPR via } G_E}
  + \underbrace{\sum_{t \in T_s} \Phi_t^{\mathrm{eff}}
      \ln \frac{\Phi_t^+}{\Phi_t^-}}_{\text{EPR via } G_s}
\end{equation}
EPR contributions from different transitions are correlated because
they share the same $G_s$-derived on/off switch ---
this is why $\mathrm{Var}[\sigma[\omega]] < \langle\sigma[\omega]\rangle$
near $N_c$ (non-MF universality class, Notes 3 and 6).

\subsection*{One-sentence summary}

\emph{$G_E$ defines the Markov generator (what states are reachable and
how fast); $G_s$ defines a dynamic filter on that generator (which
transitions are active at each state); together they produce a CTMC with
absorbing subspaces that correspond to the biological cell-fate attractors,
and whose EPR decomposes along the $G_E / G_s$ partition.}

% ─────────────────────────────────────────────────────────────────────────────
\section*{Note 10 — ``Ratchet'' Terminology: Contestable Definition}
% ─────────────────────────────────────────────────────────────────────────────

\textbf{Date:} 2026-07-16.

\subsection*{Three distinct ``ratchet'' concepts (often conflated)}

\begin{description}
  \item[Feynman ratchet-and-pawl (1963)] Mechanical device with a pawl.
    Feynman proved: at thermal equilibrium this cannot produce work.
    Only operates with a temperature gradient. \emph{Lesson: asymmetry
    alone is insufficient; non-equilibrium driving is essential.}
  \item[Brownian / Flashing ratchet (Ajdari \& Prost 1992, Magnasco 1993)]
    Spatially periodic asymmetric potential that switches on/off.
    Directed motion emerges from \emph{rectification of thermal
    fluctuations} via spatial asymmetry.
    Canonical modern ratchet.
  \item[Molecular motor ratchet (Astumian 1994+)] Kinesin, myosin.
    ATP hydrolysis creates asymmetric free energy landscape; thermal
    fluctuations are rectified into directed stepping.
    Genuine Brownian ratchet driven by chemical energy.
\end{description}

\subsection*{Why ``entropic ratchet'' for POSet cascade is contestable}

\begin{center}
\begin{tabular}{p{4cm}p{4cm}p{4cm}}
\toprule
Property & True Brownian ratchet & POSet cascade \\
\midrule
Source of directionality & Spatial asymmetry of fluctuating potential
  & Accumulated affinities $\sum \Delta\mu_k$ \\
Role of thermal fluctuations & Actively rectified (the mechanism)
  & Present but not the primary mechanism \\
Reverse steps & Allowed, probabilistically suppressed
  & Forbidden structurally (PreemptionCheck) or by large $\Delta\mu$ \\
Non-equilibrium driving & Required & ATP hydrolysis at each level \\
\bottomrule
\end{tabular}
\end{center}

In a Brownian ratchet, thermal fluctuations \emph{are} the mechanism.
In the POSet cascade, the directionality comes from the free energy
gradient; fluctuations make the process stochastic rather than
deterministic.

\subsection*{More precise description}

The POSet cascade is better described as a
\textbf{free energy cascade with absorbing barriers}: a driven directed
walk where each step is downhill in free energy and PreemptionCheck
creates absorbing states.
This is closer to ``directed walk with kinetic irreversibility'' than
to a ratchet proper.

\subsection*{Where ratchet concept converges near $N_c$}

Near $N_c$, the commitment moment IS ratchet-like: a stochastic SinI
burst (genuine thermal fluctuation) is rectified into an irreversible
fate decision by the asymmetric free energy landscape of the downstream
execution programme.
At this specific moment --- the threshold crossing --- the system
behaves as a Brownian ratchet: fluctuation $\to$ rectification $\to$
directed transition.
The ratchet concept is most precise at the commitment moment, not for
the cascade as a whole.

\subsection*{Manuscript clarification (Theory II \S4.3 footnote)}

\begin{quote}
\emph{``We use `entropic ratchet' in the loose sense of a directed
free-energy-consuming cascade, distinct from a Brownian ratchet
(Ajdari \& Prost 1992, Magnasco 1993), where spatial asymmetry
rectifies thermal fluctuations into directed motion.
The directionality in the POSet cascade arises from accumulated
chemical potential differences, not from spatial asymmetry of a
fluctuating potential.
The ratchet concept is most precise at the commitment moment itself
(near $N_c$), where a stochastic SinI burst is rectified into an
irreversible fate decision by the asymmetric free energy landscape
of the downstream execution programme.''}
\end{quote}

% ─────────────────────────────────────────────────────────────────────────────
\section*{Note 11 — Time-Reversibility and PreemptionCheck Cascades}
% ─────────────────────────────────────────────────────────────────────────────

\textbf{Date:} 2026-07-17.

\subsection*{Definition: time-reversibility for a CTMC}

A Markov process is time-reversible iff it satisfies detailed balance:
$k(x \to y)\,P_{\mathrm{ss}}(x) = k(y \to x)\,P_{\mathrm{ss}}(y)$
for all states $x,y$.
Equivalently: $P[\omega] = P[\tilde\omega]$ for every trajectory ---
$\sigma[\omega] = 0$ everywhere.

\subsection*{PreemptionCheck creates absorbing states — inherently time-irreversible}

PreemptionCheck drives the cascade toward absorbing subspaces.
For an absorbing state $x^*$: $k(x^* \to y) = 0$ for all $y \neq x^*$,
$P_{\mathrm{ss}}(x^*) > 0$, $P_{\mathrm{ss}}(x_{\mathrm{pre}}) \to 0$.
The time-reversed rate:
\[
  \tilde k(y \to x) = k(x \to y)\,\frac{P_{\mathrm{ss}}(x)}{P_{\mathrm{ss}}(y)}
\]
For $x = x_{\mathrm{pre}}$: $P_{\mathrm{ss}}(x_{\mathrm{pre}}) \to 0$
so $\tilde k \to 0$.
\textbf{Absorbing states are structurally incompatible with
time-reversibility --- a theorem, not an approximation.}

\subsection*{Three nested levels of violation}

\begin{description}
  \item[Level 1 — Global: absorbing boundary.]
    The committed state cannot be left. The time-reversed process has
    no dynamics from the committed state. Strongest violation:
    not just $\sigma[\omega] \gg 0$ but $P[\tilde\omega] = 0$ exactly.
  \item[Level 2 — Local: detailed balance violation at each cascade level.]
    $P_{\mathrm{ss}}(M_{\mathrm{pre}}) \ll P_{\mathrm{ss}}(M_{\mathrm{post}})$
    as the system approaches the absorbing boundary.
    Detailed balance requires $k(M_{\mathrm{post}} \to M_{\mathrm{pre}})
    \approx 0$ even when PreemptionCheck would nominally allow it.
  \item[Level 3 — Structural: mathematical time-reverse $\neq$ biochemical reverse.]
    The time-reverse of ``KinA phosphorylates Spo0F'' is that reaction
    un-happening. The biochemical reverse is ``RapA dephosphorylates Spo0F''
    --- a different enzyme, different rates, different Markov process.
    PreemptionCheck captures this: the reverse pathway does not satisfy
    the same PreemptionCheck conditions as the forward path.
\end{description}

\subsection*{$N_c$ as locus of maximal quasi-reversibility}

Near $N_c$, forward and backward rates become comparable (both attractors
equally probable).
Probability currents are most balanced --- locally minimal
time-reversibility violation.
Critical slowing-down: relaxation time diverges because net drift toward
either attractor shrinks.
$\gamma_{\mathrm{sub}} = 0.150$ (slow) means the cascade builds
irreversibility gradually as $N_0$ decreases --- the system remains
quasi-reversible over a wide nutrient window.

\subsection*{Gallavotti-Cohen connection}

\[
  \frac{P(\sigma[\omega] = +A)}{P(\sigma[\omega] = -A)} = e^A
\]
For time-reversible system: ratio $= 1$.
For committed cascade: ratio $= e^{3.72\times10^7}$.
PreemptionCheck makes the ratio exactly $\infty$ for structurally
forbidden paths --- discontinuous violation, not merely a large-but-finite
one.

\subsection*{Are these good or bad findings for SHPN?}

\begin{center}
\begin{tabular}{p{5.5cm}ll}
\toprule
Property & Assessment & Reason \\
\midrule
Absorbing states & Good & Correctly models biological irreversibility \\
Three-level violation & Good & Dual (thermodynamic + specificity) enforcement \\
$N_c$ as quasi-reversibility locus & Excellent & Non-trivial emergent prediction \\
Cascade depth $\to$ $W/k_BT$ scaling & Good & Falsifiable cross-organism prediction \\
Transient EPR vs.\ NESS EPR & Honest limitation & Needs explicit scope statement \\
$P[\tilde\omega] = 0$ from PreemptionCheck & Nuanced &
  Formally correct; outside standard Seifert domain \\
\bottomrule
\end{tabular}
\end{center}

\textbf{The genuine technical concern:}
Standard stochastic thermodynamics (Seifert 2012) is developed for
\emph{NESS systems} with a well-defined stationary distribution and
circulating probability currents.
The sporulation cascade is an \emph{absorbing Markov chain} --- no
finite-time stationary distribution exists.
The EPR computed is therefore a \textbf{transient EPR} (time-averaged
over the commitment trajectory before absorption), not a true
steady-state EPR.
The value $W/k_BT = 3.72\times10^7$ contains both genuine thermodynamic
dissipation (ATP affinities) and the absorbing-boundary contribution
(architectural PreemptionCheck).
This is a scope-of-interpretation limitation, not a correctness flaw.

\textbf{Overall verdict:}
The framework is internally consistent and biologically correct.
SHPN operates at the boundary of standard stochastic thermodynamics ---
where biochemical specificity and thermodynamic driving are both essential.
That boundary is exactly where the interesting biology lives.

% ─────────────────────────────────────────────────────────────────────────────
\section*{Note 12 — Categorising SHPN Within NESM}
% ─────────────────────────────────────────────────────────────────────────────

\textbf{Date:} 2026-07-23.

\subsection*{Level 1 — Computational substrate for single-cell NESM}

SHPN's stochastic transitions are $\tau$-leaping approximations of the
Master equation of a CTMC.
This makes SHPN the exact substrate required by Seifert's stochastic
thermodynamics: trajectory EPR
$\sigma[\omega] = \ln P[\omega]/P[\tilde\omega]$ is computable from firing
counts, and NESS characterisation follows without additional assumptions.
SHPN is not itself a thermodynamic theory; it is the \textbf{computational
machine through which thermodynamic quantities are derived from biological
topology}.

\subsection*{Level 2 — Thermodynamically typed arc system}

Generic CTMCs do not distinguish the thermodynamic role of each reaction.
SHPN's arc type system provides this:

\begin{center}
\begin{tabular}{p{2.5cm}p{1.8cm}p{2cm}p{4.5cm}}
\toprule
Arc type & $\Delta M$ & EPR & Thermodynamic role \\
\midrule
Normal & $W \neq 0$ & Non-zero & Stoichiometric flows \\
Signal-flow & $W_s \neq 0$ & Non-zero & Regulatory dissipation channels \\
Test & $0$ & Zero & Pure sensing, no thermodynamic cost \\
Inhibitor & $0$ & Zero & Structural gate, no mass coupling \\
\bottomrule
\end{tabular}
\end{center}

This makes SHPN a \textbf{CTMC with thermodynamically annotated transitions}
--- enabling the per-arc EPR decomposition that a bare CTMC cannot produce.

\subsection*{Level 3 — Subsystem boundary specification language}

In NESM of open systems (Esposito \& Van den Broeck 2010), defining the
thermodynamic boundary is the critical challenge.
SHPN provides this via the $G_E/G_s$ split:
\begin{itemize}
  \item $G_s$ arcs (signal-flow) declare the commitment subnetwork ---
    the thermodynamic ``inside''
  \item $\Phi$ (rate functions, remote sensing) specifies coupling to the
    environment --- the thermodynamic bath (``outside'')
  \item PreemptionCheck defines the cascade structure creating the absorbing
    boundary
\end{itemize}
The Markovian reservoir assumption ($\sigma_{\rm interface} = 0$) is made
precise: the bath enters only through $\Phi$, never through $G_s$ arcs.
This turns ``treat the rest of the cell as a bath'' into a formally
verifiable architectural property.

\subsection*{Level 4 — Theorem-generator for NESM structural properties}

\begin{center}
\begin{tabular}{p{4cm}p{7cm}}
\toprule
SHPN theorem & NESM consequence \\
\midrule
Acyclicity ($G_s$ is a DAG) & Thermodynamic irreversibility: no reverse
  trajectory reachable via $G_s$ \\
Hierarchical preemption & Correlated EPR silence: PreemptionCheck creates
  collective on/off switching $\Rightarrow$ Var[$\sigma$] suppressed
  $\Rightarrow \gamma < 0.5$ \\
Basin partition & Absorbing states: committed CTMC has
  $P(\text{recommit}) = 0$ \\
Threshold computability & $M_{\rm commit} = \theta + W_s$ gives the
  structural commitment energy barrier from topology \\
\bottomrule
\end{tabular}
\end{center}

\subsection*{Precise NESM categorisation}

\begin{quote}
\emph{SHPN is a typed stochastic Petri net whose type system implements
the thermodynamic role classification of stochastic thermodynamics:
stoichiometric mass flows (normal arcs), regulatory dissipation channels
(signal-flow arcs), zero-EPR sensing elements (test arcs), and a cascade
ordering (POSet $\lambda$ + PreemptionCheck) that structurally enforces
the thermodynamic irreversibility of commitment.}
\end{quote}

More compactly: \textbf{SHPN is the formal language in which the stochastic
thermodynamics of biological commitment is expressible and computationally
tractable.}

\subsection*{Why this matters for the two papers}

\begin{itemize}
  \item The \textbf{PLOS ONE paper} establishes SHPN as a formal language
    (syntax + semantics).
  \item The \textbf{NESM paper} establishes SHPN as a thermodynamic substrate
    (physical interpretation).
\end{itemize}

The two papers are formally dual: one proves what the formalism CAN express;
the other proves what the formalism CANNOT HELP BUT express thermodynamically.
The formal properties of the arc type system (proved in PLOS ONE) are the
structural basis for the thermodynamic properties (derived in the NESM paper).

% ─────────────────────────────────────────────────────────────────────────────
\section*{Note 13 — $C$ as a Thermodynamic Boundary Operator}
\label{note:C_thermodynamic}
% ─────────────────────────────────────────────────────────────────────────────

The compartment assignment $C: P \to \text{Compartments}$ is typically
described as a simulation-level parameter that converts token counts to
concentrations.  This note argues it is also the \emph{thermodynamic boundary
operator} of the formalism --- the element of the 13-tuple that determines
where entropy is produced, how much dissipative work each firing event costs,
and where basin boundaries are located in the free-energy landscape.

\subsection*{Chemical potential through concentration scaling}

The chemical potential of species $p$ in compartment $C(p)$ with volume
$V_{C(p)}$ is
\begin{equation}
  \mu_p = \mu_p^\circ + k_BT \ln\!\left(\frac{M(p)}{N_A \cdot V_{C(p)}}\right).
  \label{eq:note13_chempot}
\end{equation}
The same token count $M(p)$ produces different $\mu_p$ depending on which
compartment $C$ assigns $p$ to.  $C$ therefore controls the thermodynamic
driving force of every transition that consumes $p$ --- it is not neutral
bookkeeping but a physical input to the propensity functions.

\subsection*{Local detailed balance across compartment boundaries}

For a transition $t$ moving mass from compartment $C_1$ to compartment $C_2$,
local detailed balance requires
\begin{equation}
  \frac{k_t^+}{k_t^-} = \exp\!\left(-\frac{\Delta G_t}{k_BT}\right),
\end{equation}
where $\Delta G_t = \sum_i \nu_i \mu_i$ depends on the volumes assigned by
$C$.  A transition spanning two compartments ($C(p_{\mathrm{in}}) \neq
C(p_{\mathrm{out}})$) has $\Delta G_t \neq 0$ even at equal token counts ---
compartment boundaries defined by $C$ are the \emph{structural origin of
irreversibility}.

\subsection*{Entropy production rate is compartment-sensitive}

The EPR for transition $t$:
\begin{equation}
  \dot\sigma_t = (k_t^+ - k_t^-)\ln\frac{k_t^+}{k_t^-} \geq 0.
\end{equation}
Since $k_t^+/k_t^- = e^{-\Delta G_t / k_BT}$ and $\Delta G_t$ depends on
$V_{C(p)}$ via \cref{eq:note13_chempot}, the EPR is zero only when all
connected species are in the same compartment at equilibrium concentrations.
A transition that spans a compartment boundary is a \emph{non-equilibrium
steady-state driver by construction}.

\subsection*{The B. subtilis case: septation changes $C$ at runtime}

The septation event ($T_{\text{septation}}$) is the transition that changes
$C$ for a subset of places at runtime --- it creates the forespore / mother-cell
compartment split.  Before septation, $\sigma^F$ and $\sigma^E$ share a common
volume; after, they are thermodynamically isolated.
\begin{itemize}
  \item \textbf{Formalism:} $C$ changes its mapping for forespore places
    post-septation.
  \item \textbf{Thermodynamics:} the free-energy cost of maintaining the
    post-septation asymmetry appears as sustained EPR at the septum.
  \item \textbf{NESM:} the Waddington basin boundary $M_{\text{commit}}$ is
    exactly the token count at which $\Delta G_{\text{septation}} < 0$ becomes
    irreversible --- the curvature of the free-energy landscape changes from
    double-well to single-well at precisely this point.
\end{itemize}

\subsection*{Connection to the POSet ratchet (Note 10)}

The layer function $\lambda$ defines the \emph{direction} of the ratchet
(the hierarchy of commitment decisions).  $C$ defines the \emph{thermodynamic
cost} of each ratchet click: the free-energy drop $-\Delta G_t$ per firing
through a compartment boundary is the energy dissipated per commitment event.
The ratchet is not merely structural (acyclic $G_s$) --- it is thermodynamically
powered by the chemical potential differences that $C$ establishes.

\subsection*{One-sentence summary}

$C$ is the formal bridge between Petri net topology and stochastic
thermodynamics: it converts token counts into chemical potentials, making
compartment boundaries the structural origin of entropy production and
irreversibility in the basin partition.

% ─────────────────────────────────────────────────────────────────────────────
\section*{Note 14 — System, Bath, and Observer Definition in Single-Cell NESM}
\label{note:system_boundary}
% ─────────────────────────────────────────────────────────────────────────────

In classical NESM the analysis requires three things declared upfront:
\emph{system} (degrees of freedom tracked explicitly), \emph{bath}
(degrees of freedom marginalized), and \emph{observer} (what can be
measured).  For a single cell none of these are unique --- they depend on the
question.  Life solved this problem by constructing \emph{nested}
thermodynamic boundaries (plasma membrane, nuclear envelope, mitochondrial
inner membrane, septum).  This note develops how SHPN encodes all four levels
of boundary and what NESM restraints each level imposes.

\subsection*{The four-level boundary hierarchy}

\paragraph{Level 1 --- Thermal boundary (implicit).}
\textbf{Bath:} water and unmodeled vibrational / rotational degrees of freedom
at $T \approx 310\,\mathrm{K}$.
\textbf{System:} all places $P$.
\textbf{Boundary:} the modeller's coarse-graining decision --- every molecular
species \emph{not} in $P$ is marginalized into the thermal bath.
The observer \emph{is} the model: it tracks exactly $M(p)$ for $p \in P$
and ignores everything else.

\noindent Restraint: EPR computed from $P$ is a \emph{lower bound} on the
true cellular EPR.  Hidden variables (unmodelled species) contribute hidden
entropy production.  Every place omitted from $P$ moves dissipation into the
invisible bath.

\paragraph{Level 2 --- Chemical boundary (source/sink transitions).}
\textbf{Bath:} chemical reservoirs --- external nutrients, ATP regeneration,
waste sinks.
\textbf{Boundary:} transitions flagged \texttt{is\_source} / \texttt{is\_sink}.
These are the \emph{ports} through which the cell exchanges matter and free
energy with the environment.  Define:
\[
  \partial_{\mathrm{chem}}
    = \{(t,p) \in F : t \text{ is source or sink}\}.
\]
The steady-state EPR of the whole cell equals the total chemical potential
drop across $\partial_{\mathrm{chem}}$:
\[
  \dot\sigma_{\mathrm{cell}}
    = \sum_{(t,p)\in\partial_{\mathrm{chem}}}
        J_{tp}\,\frac{\mu_p^{\mathrm{in}} - \mu_p^{\mathrm{out}}}{T}.
\]
ATP and GTP pools with large, slowly changing markings function as
\emph{chemical baths}: kinetically decoupled reservoirs that maintain a
fixed chemical potential for the rest of the network.

\noindent Restraint: every SHPN model intended for NESM analysis must declare
its chemical bath explicitly via source/sink transitions.  Without them the
steady-state EPR is undefined (open system with no flux boundary).

\paragraph{Level 3 --- Informational boundary (signal places $\Psi$ and $F_s$).}
\textbf{Observer:} any agent that can detect the firing of a signal-flow arc.
\textbf{Boundary:} $F_s$ arcs --- the only events where regulatory information
is irreversibly consumed.
\textbf{System (informational):} the signal graph $G_s = (\Psi, F_s)$.

A downstream transition $t_j$ ``observes'' the state of its signal input by
checking $M(p_s) \geq \theta + W_s$ --- a \emph{binary measurement} of a
continuous quantity.  By Landauer's principle the minimum thermodynamic cost
of this measurement-and-reset is
\[
  \Delta F_{\mathrm{Landauer}} \geq k_BT \ln 2
  \quad \text{per binary gate per firing.}
\]
Every layer crossing in $G_s$ is a measurement event with an explicit
erasure cost encoded in the signal arc consumption $W_s$.
PreemptionCheck is therefore a \emph{Maxwell's demon cascade}: each layer
gate measures a lower-layer signal, makes a binary decision, and dissipates
at least $k_BT \ln 2$ per decision.

\noindent Restraint: the number of layers $\max\lambda$ is a lower bound on
the information-theoretic cost (in units of $k_BT$) of the full commitment
decision.  Models that collapse all commitment into a single transition
underestimate this cost by $(\max\lambda - 1) \times k_BT \ln 2$.

\paragraph{Level 4 --- Spatial/compartmental boundary ($C$ partition).}
\textbf{Subsystem $\alpha$:} all places $p$ with $C(p) = \alpha$, volume
$V_\alpha$.
\textbf{Boundary:} transitions $t$ where at least one input and one output
belong to different compartments.
\textbf{Local EPR at boundary $\alpha|\beta$:}
\[
  \dot\sigma_{\alpha|\beta}
    = J_{\alpha\beta}\,\frac{\mu_\alpha - \mu_\beta}{T} \geq 0.
\]
The cell maintains $\mu_{\mathrm{nucleus}} \neq \mu_{\mathrm{cytoplasm}}$
by spending EPR at all nuclear-pore transitions.
$C$ is the formal declaration: \emph{``I am spending free energy to maintain
this gradient.''}

\noindent Restraint: merging two biologically distinct compartments into one
in $C$ underestimates the EPR of all cross-boundary transport and wrongly
treats the combined pool as an equilibrated reservoir.

\subsection*{Summary table}

\begin{center}
\small
\begin{tabular}{lllll}
\toprule
\textbf{Level} & \textbf{SHPN element} & \textbf{System} & \textbf{Bath}
  & \textbf{EPR location} \\
\midrule
Thermal      & All $P$, implicit solvent & $P$             & Water, $k_BT$
  & Unobservable (lower bd.) \\
Chemical     & Source/sink transitions   & Interior $P$    & ATP, nutrients
  & All source/sink arcs \\
Informational & $\Psi$, $F_s$           & $G_s$           & $G_E$ continuum
  & Each $F_s$ firing $\geq k_BT\ln2$ \\
Spatial      & $C$ partition            & Compartment $\alpha$ & Compartment $\beta$
  & Cross-$C$ transitions \\
\bottomrule
\end{tabular}
\end{center}

\subsection*{The septation event as a boundary-creation event}

$T_{\text{septation}}$ is unique among transitions: it does not merely consume
tokens --- it \emph{changes the topology of the boundary set}.  After firing,
$C$ maps forespore and mother-cell places to different volumes, creating two
Level-4 NESM subsystems from one.  The EPR spike observed at commitment
(Results II) is partly the \emph{thermodynamic cost of boundary creation} ---
the work of sealing the membrane and establishing the cross-septum chemical
potential gradient that drives asymmetric sigma-factor activation.

\subsection*{Four NESM restraints on SHPN model construction}

\begin{enumerate}
  \item \textbf{Declare the bath explicitly.}  Source/sink transitions +
    parameter places $\square$ define the environment.  Without them EPR is
    uncomputable.
  \item \textbf{Ground $C$ in biology.}  Each compartment must correspond to
    a real membrane-bounded volume.  Merging compartments underestimates
    cross-boundary EPR.
  \item \textbf{Count the Landauer cost of layers.}  Each layer in $\lambda$
    is a binary decision gate with minimum cost $k_BT\ln2$.  The layered
    cascade is thermodynamically expensive by design --- that cost is what
    makes the commitment irreversible.
  \item \textbf{Declare the observer coarse-graining.}  The set $P$ is the
    observer's accessible information.  EPR from the model is conditioned on
    this observer.  A finer-grained observer (larger $P$) would see higher
    EPR.  Results must be reported with $P$ specified.
\end{enumerate}

\subsection*{Why compartmentalisation is thermodynamically rational}

Life chose compartmentalisation because it enables:
(a) different chemical potentials in different volumes --- mechanical work
extraction from concentration gradients;
(b) local control of reaction environments --- kinetic specialisation without
crosstalk;
(c) information isolation --- noise is contained within a compartment, not
broadcast to the whole cell;
(d) thermodynamic separation --- each compartment can approach its own local
quasi-equilibrium between commitment events, reducing the EPR required to
maintain the resting state.

SHPN's $C$ function is therefore not a simulation convenience.  It is the
formal statement that the cell has \emph{chosen} to pay the thermodynamic
cost of maintaining multiple distinct chemical potentials simultaneously ---
and that the resulting gradients are the engine that drives irreversible
commitment.

% ─────────────────────────────────────────────────────────────────────────────
\section*{Note 15 — Stochastic Thermodynamics Gives Equalities, Not Just Directions}
\label{note:equalities}
% ─────────────────────────────────────────────────────────────────────────────

A common misconception: stochastic thermodynamics (ST) is just the second law
applied to small systems --- a collection of inequalities saying which way
processes go.  This is wrong.  ST converts inequalities into exact identities
and provides fully quantitative, trajectory-level measures.  This note
catalogues what ST actually gives, and what that means for the five
topological-thermodynamic quantities in the NESM manuscript.

\subsection*{Classical thermodynamics: directions only}

Classical macroscopic thermodynamics gives inequalities:
\[
  \Delta S \geq 0, \qquad \Delta G \leq 0.
\]
These tell you which direction a process goes --- nothing about how fast, how
far from equilibrium, or what the fluctuations look like.  The second law is
an ensemble statement; it says nothing about any individual trajectory.

\subsection*{Stochastic thermodynamics: exact equalities}

\paragraph{1. Entropy production rate --- exact number.}
\[
  \dot\sigma = \sum_\rho J_\rho \ln\frac{k_\rho^+}{k_\rho^-}
  \quad [k_B\,\mathrm{s}^{-1}].
\]
Computable directly from fluxes and rate ratios.  A specific number per
transition per instant, not a directional inequality.

\paragraph{2. Jarzynski equality --- the second law becomes an identity.}
Classical: $\langle W \rangle \geq \Delta F$ (direction only).
Stochastic:
\[
  \bigl\langle e^{-W/k_BT} \bigr\rangle = e^{-\Delta F / k_BT} \quad \text{(exact).}
\]
The exact free energy difference is recoverable from non-equilibrium work
measurements.  An inequality becomes an equality; a direction becomes a
number.

\paragraph{3. Crooks fluctuation theorem --- quantitative ratio.}
\[
  \frac{P_F(W)}{P_R(-W)} = e^{(W - \Delta F)/k_BT}.
\]
The crossing point of the forward and reverse work distributions locates
$\Delta F$ exactly.  Quantitative, not directional.

\paragraph{4. Gallavotti--Cohen symmetry --- probability of violations.}
\[
  \frac{P(\sigma = +A)}{P(\sigma = -A)} = e^{A/k_B}.
\]
Not just ``second-law violations are rare'' --- exactly how rare, as a
function of magnitude.  A hard number.

\paragraph{5. Thermodynamic uncertainty relation (TUR) --- exact precision-dissipation tradeoff.}
\[
  \frac{\mathrm{Var}(J)}{\langle J\rangle^2} \geq \frac{2k_B}{\Sigma}.
\]
To halve the noise of a molecular clock, you must quadruple the dissipation.
Not a direction --- an exact exchange rate between precision and energy cost.

\paragraph{6. Thermodynamic speed limit --- quantitative time-cost.}
\[
  \tau \geq \frac{L^2}{\dot\sigma}.
\]
Minimum time to traverse state-space distance $L$ given dissipation rate
$\dot\sigma$.  A lower bound on commitment time given available free energy.

\paragraph{7. Landauer --- exact information-work conversion.}
\[
  \Delta F_{\text{erase}} \geq k_BT \ln 2 \quad \text{per bit.}
\]
One bit of erased information costs exactly $k_BT \ln 2$ joules minimum.
The PreemptionCheck cascade has $\max\lambda$ binary gates; their total
Landauer cost is $\max\lambda \times k_BT \ln 2$, exactly.

\subsection*{Comparison table}

\begin{center}
\small
\begin{tabular}{lll}
\toprule
\textbf{Quantity} & \textbf{Classical} & \textbf{Stochastic} \\
\midrule
Direction of process         & $\Delta S \geq 0$ (ineq.) & $\dot\sigma \geq 0$ (same) \\
Rate of entropy production   & not defined               & exact: $\sum_\rho J_\rho \ln(k^+/k^-)$ \\
Free energy from non-eq.     & $\langle W\rangle \geq \Delta F$ & exact: Jarzynski equality \\
Fluctuations                 & averaged away             & exact distributions via FTs \\
Time asymmetry               & direction only            & exact ratio: Gallavotti--Cohen \\
Precision vs cost            & no statement              & exact bound: TUR \\
Speed vs cost                & no statement              & exact bound: speed limit \\
Information cost             & direction (Landauer ineq.)& exact: $k_BT\ln 2$ per bit \\
Single-trajectory EPR        & not defined               & $\sigma[\omega] = \ln P[\omega]/P[\tilde\omega]$ \\
\bottomrule
\end{tabular}
\end{center}

\subsection*{What this means for the five quantities}

All five topological-thermodynamic quantities in the NESM manuscript are
\emph{quantitative} in the ST sense:
\begin{itemize}
  \item $\eta_L$ --- dimensionless number (arc density of commitment layer).
  \item $W/k_BT$ --- how many thermal energy units the commitment dissipates;
    computable via Jarzynski from replicate work distributions.
  \item $N_c$ --- exact concentration in $\mu$M where $P(\text{sporulate})=0.5$;
    a phase-transition order parameter, not a directional bound.
  \item $\gamma_{\mathrm{sub}}$ --- a critical exponent with mean-field
    prediction $0.5$; deviations are quantitative evidence of long-range
    correlations.
  \item EPR decomposition --- $k_B\,\mathrm{s}^{-1}$ per transition per
    replicate; the budget of dissipation assigned to each biochemical event.
\end{itemize}

\subsection*{The key sentence for the manuscript}

The second law becomes an equality at the trajectory level:
$\langle e^{-\sigma[\omega]} \rangle = 1$ (integral fluctuation theorem).
The ensemble average gives the direction ($\langle\sigma\rangle \geq 0$);
the distribution of $\sigma[\omega]$ across replicates gives the Gallavotti--Cohen
ratio (the quantity).  SHPN makes both computable from the same $\tau$-leaping
output.

% ─────────────────────────────────────────────────────────────────────────────
\section*{Note 16 — Structural vs Thermodynamic Irreversibility: Non-Equivalence}
\label{note:structural_vs_thermal}
% ─────────────────────────────────────────────────────────────────────────────

\textbf{Question:} If we remove the PreemptionCheck/POSet mechanism from
SHPN, do the irreversibility theorems of stochastic thermodynamics guarantee
the same results?

\textbf{Answer: No.}  They guarantee different things and are not equivalent.

\subsection*{What ST irreversibility guarantees}

The Gallavotti--Cohen theorem:
\[
  \frac{P(\sigma = +A)}{P(\sigma = -A)} = e^{A/k_B}.
\]
The integral fluctuation theorem:
\[
  \bigl\langle e^{-\sigma[\omega]} \bigr\rangle = 1.
\]
Both say reverse trajectories (de-commitment) are \emph{exponentially
suppressed} but never forbidden.  The probability of de-commitment is:
\[
  P(\text{de-commit}) \sim e^{-\sigma_{\text{commit}}/k_B} > 0 \quad \text{always.}
\]
For large $\sigma_{\text{commit}}$ this is astronomically small ---
but mathematically non-zero at any finite temperature.

\subsection*{What PreemptionCheck/POSet guarantees}

Theorem~3 (Basin Partition): once below $M_{\text{commit}}$ with no regenerating
$F_s$ arc, return to $B_{\text{above}}$ has probability \textbf{exactly 0} ---
not exponentially small.  The structural arc is absent; the reverse transition
cannot fire because it does not exist.

\subsection*{The four non-equivalences}

\paragraph{1. Probabilistic vs absolute.}
ST: $P(\text{reverse}) = \varepsilon > 0$ (exponentially small). \\
SHPN: $P(\text{reverse}) = 0$ (topologically forbidden). \\
These coincide only in the limit $\sigma \to \infty$.  At any finite
temperature and finite dissipation, ST irreversibility is approximate;
structural irreversibility is exact.

\paragraph{2. Requires equilibrium reference vs topology-only.}
ST irreversibility is defined \emph{relative to} an equilibrium state:
you need $k_t^-$ (the reverse rate) to compute $\sigma$.  If reverse rates
are absent by topology, ST gives $\sigma = \infty$ (degenerate case).
SHPN irreversibility requires only: no $F_s$ arc regenerates $p_s$.
Readable from topology alone, no rate constants needed.

\paragraph{3. Temperature-dependent vs temperature-independent.}
At higher $T$, thermal fluctuations increase, EPR barriers shrink, and
$P(\text{de-commit}) \sim e^{-\sigma/k_B}$ grows.
A cell under heat stress has a less reliable ST commitment.
PreemptionCheck is temperature-blind: the arc is absent or present
regardless of $T$.  Structural enforcement is more robust under
stress conditions (heat shock, ATP depletion) where $\Delta G / k_BT$ shrinks.

\paragraph{4. Timescale agnosticism.}
ST gives the probability of de-commitment; Kramers escape gives the time:
\[
  \tau_{\text{escape}} \sim \tau_0 \, e^{\Delta G / k_BT}.
\]
For $\Delta G = 50\,k_BT$: $\tau_{\text{escape}} \sim e^{50} \approx 10^{21}$
attempt periods --- safe biologically.  But near $N_c$, $\Delta G \to 0$;
Kramers escape becomes fast and ST irreversibility collapses on biological
timescales.  PreemptionCheck still holds.

\subsection*{Set-theoretic relationship}

\[
  \text{SHPN structural irreversibility} \;\subsetneq\; \text{ST irreversibility}.
\]
Every structurally irreversible system is also thermodynamically irreversible
(positive EPR).  The converse fails: a system with large EPR but no topological
gate can be de-committed by a sufficiently large fluctuation.

\subsection*{The biological design logic}

The cell uses both mechanisms in distinct phases:

\begin{center}
\small
\begin{tabular}{lll}
\toprule
\textbf{Phase} & \textbf{Dominant mechanism} & \textbf{Role} \\
\midrule
$N_0 > N_c$ (vegetative)   & ST (high EPR, NESS)          & Maintain vegetative state \\
Near $N_c$ (decision zone) & Neither (low EPR, bistable)  & Bet-hedging: stochastic probing \\
After crossing $M_{\text{commit}}$ & SHPN structural (PreemptionCheck) & Lock decision permanently \\
Layer B (execution)        & ST (ODE-limit EPR)           & Power gene expression \\
\bottomrule
\end{tabular}
\end{center}

Near $N_c$ the EPR is low and ST alone cannot guarantee commitment permanence
--- the cell IS near-reversible by design (bet-hedging).  The structural gate
fires AFTER the fluctuation crosses the threshold, locking in the decision.

\subsection*{One-sentence answer}

ST guarantees commitment is unlikely to reverse (exponential suppression);
SHPN guarantees it cannot reverse (topological impossibility).  They coincide
only at infinite dissipation.  The cell needs both: ST for the decision cost,
PreemptionCheck for the decision lock.

\end{document}
